\section{Limitations and conclusion}
\label{sec:conclusion}

\paragraph{Limitations.} The agents are single-layer perceptrons over a fixed
shared embedding; the classes are binary, balanced and immutable; every pair may
interact, so there is no network structure; the population is closed and of fixed
size. Issues are isotropic synthetic vectors rather than embeddings of anything.
Because the dynamics anneals rather than converging, every regime we report is a
regime \emph{at a stated interaction count}, and Appendix~\ref{app:simulation}
shows how each observable moves with that count. ``Trust'' is a model construct
--- an inferred channel-noise level --- and not a measurement of anything
psychological. We have not validated any of this against behavioural or societal
data, and it would not license conclusions about human groups if we had.

\paragraph{Conclusion.} A task-irrelevant feature can do more than bias a single
prediction. When learners also infer whom to trust, a class-correlated
representation error changes the information flow of the entire population: it
breaks a symmetry of the local dynamics, and the break is amplified into
group-aligned collective structure that no individual agent chose and that the
agenda's complexity can reorder in time. Two consequences seem worth carrying
away. The first is that the vulnerability comes from the very machinery that
makes each agent a good inferential citizen --- drawing conclusions about
\emph{sources} as well as about the world --- so giving interacting agents a
well-founded model of one another's reliability is not a neutral improvement.
The second is that the label starts out carrying no information and does not stay
that way: once opinions align with class, class genuinely predicts opinion, and
an agent audited against the realised distribution would find nothing to correct.
Population-level measurement is not a refinement of per-agent evaluation here; it
is the only thing that can see the effect at all.
